\cotefile{157-4}
\scanpage{148}
\mspage{22}
% archive sheet number: 147

de cardinal majoré par
\[
\underbrace{r^{\pi_0}}_{\substack{\text{nb de \unsure{classes|ces}}\\\text{syst. inductifs}\\\text{dans }\underline M^{\pi},\text{ à}\\\text{isom près}}}
\mathbin{\cdot}
\underbrace{s^{r^{\pi_0}}}_{\substack{\text{nb des}\\\text{flèches d'un}\\\text{tel syst. ind.}\\\text{dans }X}}
= s^{r^{\pi_0}}.
\]

Mais il faut maintenant faire varier $I$,
et nous \unsure{porter|plonger} dans la cat des ens. ordonnés
(filtrants etc) de cardinal $\inserted{\lambda=}\deleted{\unsure{I}}r^{\pi_0}\deleted{\unsure{\longrightarrow}}$.
On voit de suite que cet ensemble de classes d'isomorphie (qui peuvent s'interpréter comme un ens.
quotient de \deleted{\unsure{l'ens}} \inserted{un certain} ens. de parties d'un ens.
fixe $I_0$ de card $\lambda$) est de cardinal $\leq 2^\lambda$.
On trouve donc finalement
\[
\operatorname{card}\widetilde{M^r_{/X}}\leq s^{(r^{\pi_0})}\mathbin{\cdot}2^{(r^{\pi_0})}=s^{(r^{\pi_0})},
\]
sous condition que
\[
r\geq\operatorname{Sup}(\pi_1,\unsure{\mathfrak C_0|C_0}),\qquad s\geq\operatorname{Sup}(\unsure{\pi_1|\mathfrak C_1},\mathfrak C_1).
\]

Il faut prouver l'inégalité similaire pour
$\operatorname{card}\widetilde{\underline{\mathrm{Fl}}^{\,r}M}$ (3.9.). Or, \deleted{\unsure{étant entendu}} \inserted{comme}
pour \deleted{\unsure{x}} $A,B$ dans $M^r$, on a si $r\geq\operatorname{Sup}(\pi_1,C_1)$
\[
\operatorname{card}\operatorname{Hom}_{M}(A,B)\leq r^{(r^{\pi_0})}
\]
en vertu de (3.8) où on fait $s=r$, on trouve que
\[
\begin{aligned}
\operatorname{card}\widetilde{\underline{\mathrm{Fl}}^{\unsure{r}}(M^r_{/X})}
\deleted{\underbrace{\vphantom{\mathrm{Fl}}}_{\unsure{\mathrm{Fl}(M^r)}}}
&\leq\left(\operatorname{card}\widetilde{\deleted{\unsure{M}}M^r_{/X}}\right)^2\mathbin{\cdot}r^{(r^{\pi_0})}\\
&\leq s^{(r^{\pi_0})}\mathbin{\cdot}r^{(r^{\pi_0})}=s^{(r^{\pi_0})}\\[-2pt]
&\hspace{6em}\text{si }s\geq r.
\end{aligned}
\]
% slanted
\margin{cf \unsure{corol.|contr.} \unsure{pge|pg} 23}
(On va regarder plus bas le cas $s\leq r\unsure{{}^{*}|{}^{+}}$\unsure{.}). C'est bien
la formule (3.9.)
\scanpage{149}
\mspage{23}
% archive sheet number: 148

Venons-en à la catégorie $M^r$ elle $=$. L'argument
de la page 21 montre que l'ens des \inserted{classes d'isom.}
objets de $M^r$ indexés par un $I$ fixé, de cardinal $\leq r^{\pi_0}$, comme
$\deleted{\unsure{X}}A=\underset{i}{\underline{\lim}}\,A_i$ avec $A_i\in\unsure{\mathrm{Filt}}^{\pi_0}$, est de cardinal
majoré par $r^{\pi_0}$ (si $r\geq C_0$). Si on fait varier
$I$, on trouve \deleted{\unsure{$r^{\pi}$}} une majoration par $r^{\pi_0}\mathbin{\cdot}2^{(r^{\pi_0})}
=2^{(r^{\pi_0})}$. On en conclut comme tantôt que
l'ens $\widetilde{\mathrm{Fl}\, M^r}$ est de cardinal majoré par
\[
\left(2^{(r^{\pi_0})}\right)^2\times r^{(r^{\pi_0})}
\qquad(\text{si }r\geq\deleted{\unsure{\pi}}\operatorname{Sup}(\pi_1,C_1)),
\]
donc par $r^{(r^{\pi_0})}$. Je dis qu'on a
\[
r^{(r^{\pi_0})}=2^{(r^{\pi_0})},
\]
en effet, on a $\geq$, et pour prouver $\leq$, on
note que $r<2^r$, donc $r^{(r^{\pi_0})}\leq(2^r)^{(r^{\pi_0})}=2^{r\mathbin{\cdot}r^{\pi_0}}=2^{r^{\pi_0}}$,
OK ! Cela établit donc 3.7. \deleted{On en déduit} \inserted{Cet argument}
\unsure{donne} en fait le

\uline{Lemme 3.\unsure{?|2}.} Soient $r$ un cardinal, $r'$ un cardinal $\geq r$. Alors
on a
\[
r^{r'}=2^{r'}\qquad(=s^{r'}\text{ pour tout }2\leq s\leq r)
\]

\uline{Corollaire 1} Soient \deleted{$r\geq s\geq\operatorname{Sup}(\pi_1,C_1)$.} $r\geq\operatorname{Sup}(\pi_1,C_1)$,
et $A,X\in\unsure{\mathrm{Pfilt}|\mathrm{Filt}}^{\unsure{\pi|r}}$. Alors
\[
\operatorname{Card}\operatorname{Hom}_M(A,X)\leq 2^{(r^{\pi_0})}
\qquad(=2^r\text{ si }r=2^c,\ \inserted{\text{avec }}c\geq\pi_0)
\tag*{(3.10.),}
\]
On applique 3.8., \deleted{\unsure{avec}} qui donne la majoration
$r^{(r^{\pi_0})}$, et \unsure{le|les} lemme 3.2.

\deleted{\uline{Corollaire 2}} On notera qu'on trouve la même
majoration que pour le cardinal de $\widetilde{\underline{\mathrm{Fl}}(M^r)}$
\unsure{dont mention} (3.7) ! De $=$ (3.9.) donne

\uline{Corollaire 2} Soit $r\geq\operatorname{Sup}(\pi_1,C_1)$, et $\unsure{X|\alpha}\in\operatorname{Ob}M^r$. Alors
\[
\operatorname{Card}\widetilde{\underline{\mathrm{Fl}}(M^r_{/X})}\leq 2^{(\unsure{r}^{\pi_0})}
\qquad(=2^r\text{ si }r=2^c,\text{ avec }c\geq\pi_0).
\tag*{(3.11).}
\]
c'est donc la formule (3.9) pour $s\leq r$, lorsque on
suppose.
\scanpage{150}
\mspage{24}
% archive sheet number: 149

Revenons à la majoration cardinale des
$X_\alpha(f)$, où $f:X\unsure{\text{---}|\longrightarrow}Y$ \unsure{décrits|décrit|reste} dans $\mathrm{Fl}^{\unsure{s}}M$, i.e.
$X,Y\in\mathrm{Filt}^s$. On suppose de plus $C_0\subset\deleted{\unsure{\mathrm{Filt}}}\mathrm{Fl}^rM$,
\deleted{\unsure{et}}
\[
\left\{\begin{aligned}
&C_0\subset\mathrm{Fl}^rM=\mathrm{Fl}\,M^r,\quad
(f\unsure{:|=}X\text{---}Y)\in\mathrm{Fl}^sM\quad\text{i.e. }X,Y\in\mathrm{Filt}^s\\
&\text{avec }s\geq r
\end{aligned}\right.
\tag*{(3.10)}
\]
\deleted{\unsure{Alors} pour $C_0$, \unsure{notons qu'aux.|notons qu'en} d'objets}
\[
\deleted{\operatorname{Card}\Delta(C_0,f)\leq\operatorname{Card}M^r_{/X}\times\operatorname{Card}M^r_{/Y}}
\]

Il y a une difficulté idiote, c'est que
\deleted{\unsure{de ce}} $M^r$ est équivalente à une petite catégorie,
mais n'est pas elle-même petite. Même si
$C_0$ ne contient qu'une seule classe d'isomorphie de flèches de $M^r$, son cardinal pourrait
être (pour $r$ fixé) arbitrairement élevé. Aussi
dans ce qui suit, je \inserted{dois} supposer \deleted{que le}
\inserted{la condition} \deleted{nb d'éléments, cardinal de \unsure{la}}
de ``\unsure{non pléthoricité} cardinale'' pour $C_0$

\[
\begin{aligned}
&\forall i\in C_0,\quad\deleted{\text{l'ens des}}\inserted{\text{désignant}} \text{ par }C_0(i)\text{ l'ens des}\\
&\qquad\deleted{\unsure{A}}i'\in C_0\text{ telles que }i\simeq i'\text{ dans }\mathrm{Fl}\,M,\text{ j'impose}\\
&\qquad\text{que }\operatorname{Card}C_0(i)\leq 2^{(r^{\pi_0})}
\end{aligned}
\tag*{(3.1\unsure{1|b})}
\]
(ce qui est \inserted{aussi} la \inserted{valeur de la} majoration pour le cardinal de
$\widetilde{\underline{\mathrm{Fl}}^rM}=\widetilde{\underline{\mathrm{Fl}}\,M^r}$)

Cela entraîne donc \deleted{que} l'on a
$\inserted{\operatorname{Card}C_0\leq}\deleted{\unsure{\operatorname{Card}}}(2^{(r^{\pi_0})})^2=\operatorname{Card}2^{r^{\pi_0}}$.
\[
\operatorname{Card}C_0\leq 2^{(r^{\pi_0})},
\tag*{(3.12)}
\]
qui est équivalente à (3.1\unsure{1|b}) (moyennant $C_0\subset\mathrm{Fl}^rM$,
postulée dans (3.10)).

Il s'ensuit \inserted{par (3.8)} que l'on a
\[
\operatorname{Card}\Delta(C_0,f)\leq
\underbrace{\operatorname{Card}(C_0)}_{\leq 2^{(r^{\pi_0})}}
\mathbin{\cdot}s^{(r^{\pi_0})}\mathbin{\cdot}s^{(r^{\pi_0})}\leq s^{(r^{\pi_0})}.
\]

D'autre part, les $X_D$ dont on prend la somme
amalgamée sous $X$, pour obtenir $X_1(f)$,
sont dans $M^s$, si on fait l'hypothèse suivante
\scanpage{151}
\mspage{25}
% archive sheet number: 150

\[
\left|\begin{aligned}
&\text{Les }M^\pi=\underline{\mathrm{Filt}}^\pi\text{ sont stables par sommes}\\
&\text{amalgamées simples, et par }\underrightarrow{\lim}\text{ filtrantes de cardinal}\\
&\leq\pi\text{ (pas néc. } \unsure{\text{grandes devant}|\text{grands devant}}\ \pi_0),
\end{aligned}\right.
\tag*{(3.13)}
\]
compte tenu que $r\leq s$ donc $M^r\subset M^s$, donc
$A,B,X$ dans $X_D=\deleted{\unsure{A}}B\amalg_A X$ sont dans $M^s$. Les sommes
amalgamées partielles \uline{finies} des $X_D$ sont
donc aussi dans $M^s$, par récurrence sur le
cardinal de l'ens d'indices, \deleted{\unsure{d'où}} la somme
amalgamée infinie apparaît ainsi comme
$\underrightarrow{\lim}$ filtrante d'objets de $M^s$, l'ens d'indices
étant l'ens. $\mathfrak P_f^*\deleted{\unsure{\subset}}\Delta(C_0,f)$ des parties finies
non vides de $\Delta(C_0,f)$. Cet ens. est de
cardinal majoré par \deleted{$\operatorname{Card}\Delta(C_0,f)\mathbin{\cdot}\unsure{\aleph}$} $s'$, \deleted{dès
que} si $s'$ est un cardinal \uline{infini} tel que
$\operatorname{card}\Delta(C_0,f)\leq s'$. On peut donc prendre $s'=s^{(r^{\pi_0})}$.
On a donc prouvé, compte tenu de (3.13)
appliqué à $M^{s'}$, le

\uline{Lemme} 3.3. On suppose (3.10) i.e. $C_0\subset\mathrm{Fl}^rM$,
$f\in\mathrm{Fl}^sM$, avec \deleted{$r\unsure{\leq}s$} \inserted{$s\geq r\geq\operatorname{Sup}(\pi_1,C_1)$},
et de plus que $C_0$ satisfait la condition de non pléthoricité (3.11.)
(ce qu'on peut toujours supposer, quitte à réduire
$C_0$ sans changer $\widetilde C_0=\Phi$ \deleted{\unsure{ni}} et $C_{0*}=\Psi$), et
que la filtration cardinale satisfait la condition
de stabilité (\unsure{anodine} !) (3.13.). Posons
\[
s'=s^{(r^{\pi_0})}.
\]
Alors $X_1(C_0,f)$ est dans $M^{s'}$. Plus précisément et
plus généralement, pour $C_0$ et $r$ fixés (avec $C_0\subset\mathrm{Fl}^rM$
et satisfaisant la condition de non pléthoricité
(3.11) relativement à $r$), pour \deleted{tout} cardinal
\scanpage{152}
\mspage{26}
% archive sheet number: 151

de la forme
\[
s'=\deleted{\unsure{s^{\pi}}}s^t,\ \text{avec }t\geq r^{\pi_0},\ s\geq2
\tag*{(3.15.)}
\]
(\uline{\unsure{d'où|alors}} $s'\geq2^{(r^{\pi_0})}>r^{\pi_0}\geq r$), $\underline{\mathrm{Fl}}^{\,s'}M$ est stable sous le
foncteur $p_1^{C_0}:\underline{\mathrm{Fl}}\,M\unsure{\text{---}|\longrightarrow}\underline{\mathrm{Fl}}\,M$.

En effet, si on applique la première partie du
lemme \unsure{au} couple $(\unsure{r|s},s')$, on trouve que
$X_1(f)$ est dans $M^{s''}$, avec $s''={s'}^{(r^{\pi_0})}$. Mais
\[
s''={s'}^{(r^{\pi_0})}=(s^t)^{(r^{\pi_0})}
=s^{t.r^{\pi_0}}\underset{\substack{\text{car }t\geq r^{\pi_0}\\\text{donc }tr^{\pi_0}=t}}{=}
s^t=s'
\]
donc on voit que $f\in\mathrm{Fl}^{s'}$ implique $X_1(f)\in M^{s'}$, ou
ce qui qui revient au même, $p_1(f)\in\mathrm{Fl}^{s'}$.

\uline{Corollaire} Soient $C_0,r,s'=s^t$ ($t\geq r^{\pi_0}$, \inserted{$s\geq2$}) comme
dessus, et soit $\alpha$ un nombre ordinal \deleted{de} \inserted{tel}
\deleted{cardinal} que
\[
\operatorname{card}\alpha\leq s'=s^t.
\]
Alors $\underline{\mathrm{Fl}}^{\,s'}M$ est stable par le foncteur $p_\alpha=p_\alpha^{C_0}$.

On le voit aussitôt par récurrence transfinie,
\unsure{compte} \deleted{\unsure{donc}} \unsure{tenu} du lemme 3.3, et de l'hyp. de
stabilité 3.13.

Pour résumer

\uline{Théorème 2} \unsure{Soit|Soit,} $M$ une catégorie stable par \inserted{petites} $\underrightarrow{\lim}$
filtrantes, et munie d'une filtration cardinale
$(\mathrm{Filt}^\pi)_{\pi\geq\pi_0}$, telle que les $\mathrm{Filt}^\pi$ soient stables
par $\underrightarrow{\lim}$ filtrantes de cardinal $\leq\pi$ (pas forcément
grandes devant $\pi_0$), et par \inserted{les} sommes amalgamées simples
qui existent dans $M$. On suppose que les éléments
de $\mathrm{Filt}^{\pi_0}$ soient accessibles, et on désigne par
$\pi_1$ le plus petit cardinal $\geq\pi_0$, tel que les objets de
$\mathrm{Filt}^{\pi_0}$ soient $\pi_1$-accessibles. On pose \deleted{de plus}
\scanpage{153}
\mspage{27}
% archive sheet number: 152

\unsure{en|---}
\[
C_1=\operatorname{card}\underbrace{\mathrm{Fl}\,M^{\pi_1}}=\operatorname{card}\mathrm{Fl}^{\pi_1}M
\]
(où pour tt $\pi\geq\pi_0$, $M^\pi$ désigne la sous-catégorie \inserted{str.}
pleine de $M$ définie par $\mathrm{Filt}^\pi\deleted{\unsure{s}}$, et $\mathrm{Fl}^\pi M=\mathrm{Fl}\,M^\pi$).

Soit $C_0$ une petite partie de $\mathrm{Fl}\,M$, formée de
flèches coquarrables, et posons
\[
\Phi=\widetilde C_0=\Psi^{*}_{\deleted{\unsure{*}}},\quad\text{où}\quad\Psi=C_{0*}=\Phi_*.
\]
Soit $r$ un ordinal $\geq\operatorname{Sup}(\pi_1,C_1)$, et tel que l'on
ait $C_0\subset\mathrm{Fl}^rM$. On suppose que $C_0$ est \unsure{irrédundant|irrédondant}
(i.e. deux objets de $C_0$ isomorphes sont égaux)
plus généralement, que $C_0$ satisfait la \inserted{très anodine}
condition de non pléthoricité (3.11.) relativement à $r$.
On considère la suite transfinie de foncteurs
$p_\alpha=p_\alpha^{C_0}$ sur $\underline{\mathrm{Fl}}\,M$ (cf. \unsure{n°|N.} 1), \deleted{Alors} \deleted{et on rappelle}
\deleted{\unsure{que ce qui}} donnant lieu aux factorisations
fonctorielles d'une $\underline f\in\underline{\mathrm{Fl}}\,M$, $f:X\text{---}Y$, en
\diagram{\begin{tikzcd}[ampersand replacement=\&,column sep=large]
\text{(3.16)} \& \unsure{x|X} \arrow[r,"i_\alpha(f)"] \& X_\alpha(f) \arrow[r,dash,"p_\alpha(f)","{\unsure{\text{---}|\to}}"'] \& Y .
\end{tikzcd}}

On a alors ce qui suit :

$1^\circ)$ On a (pour mémoire) $i_\alpha(f)\in\Phi$ pour tout $\alpha$, et
de plus, \uline{si $\alpha$ est grand devant $r^{\pi_0}$}, on a aussi
$p_\alpha(f)\in\Psi$. Si $r'$ est un cardinal $>r^{\pi_0}$ (p. ex. le
\deleted{\unsure{ses}} cardinal successeur de $r^{\pi_0}$), on peut prendre
pour $\alpha$ le premier ordinal dont le cardinal soit égal
à $r'$.

$2^\circ)$ \deleted{Supposons $\alpha$ choisi grand devant $r^{\pi_0}$. Alors} Pour
tout cardinal $s$ de la forme
\[
s=c^t\quad\text{avec }t\geq r^{\pi_0},\quad c\geq2
\]
(donc $s=2^t$ si on prend $c\leq t\unsure{\}}$, cf lemme 3.\unsure{2}), $\underline{\mathrm{Fl}}^{\,s}M$
est stable sous les $p_\alpha$ tels que $\operatorname{card}\alpha\leq s$. Si on
prend p. ex. pour $\alpha$ le premier ordinal dont le
cardinal soit $>r^{\pi_0}$, alors $\underline{\mathrm{Fl}}^{\,s}M$ est stable par
$p_\alpha\unsure{{}^{*}}$, qui par $1^\circ)$ réalise donc une factorisation fonctorielle
(3.\unsure{16|15}) par $i_\alpha(f)\in\Phi$, $p_\alpha(f)\in\Psi$.
\scanpage{154}
\mspage{28}
% archive sheet number: 153

\uline{Dém} Tout a été prouvé\unsure{,|\}} si on se rappelle que les
objets de $M^r$ sont $r^{\pi_0}$-accessibles (cf. 3.2).

\deleted{\uline{Corollaire} \unsure{La conclusion|Les conclusions} de stabilité \unsure{reste valable|restent valables},
lorsque $s$ est}

4. Compléments sur les structures de Quillen.

\uline{Définition 4.1} Soit $M$ une catégorie. On appelle
\uline{précouple de Quillen} dans $M$ un couple
$(\Phi,\Psi)$ de parties de $\mathrm{Fl}\,M$, satisfaisant les
deux conditions suivantes

a) On a $\Phi\longleftrightarrow\Psi$, i.e. $\Psi\subset\Phi_*$, ou ce qui revient
au $=$, $\Phi\subset\Psi^*$ \deleted{\unsure{et}} -- soit encore : Pour
tout carré commutatif dans $M$

\diagram{\begin{tikzcd}[ampersand replacement=\&,column sep=large]
\& A \arrow[r,"u"] \arrow[d,dash,"i"',"{\unsure{\text{---}|\downarrow}}"] \& X \arrow[d,"p"]\\
\text{(4.1.1.)} \& B \arrow[r,dash,"v"'] \arrow[ur,dotted,"w"] \& Y
\end{tikzcd}}
avec $i\in\Phi$, $p\in\Psi$, il existe $w$ telle que
$wi=u,\ pw=v$.
\margin{\unsure{i.e.|.2.}}

b) Toute $f\in\mathrm{Fl}\,M$ se factorise en
\[
f=pi,\quad\text{avec }i\in\Phi,\ p\in\Psi.
\]

On écrit
\[
\Phi\mathrel{\overset{\vphantom{I}}{\longleftrightarrow}\mkern-19mu\mid\mkern9mu}\Psi
\tag*{(4.1.2)}
\]
pour indiquer que $(\Phi,\Psi)$ est un \inserted{pré-}couple
de Quillen.

On \deleted{appelle} dit que le couple de Quillen $(\Phi,\Psi)$
est \uline{clos}, si $\Phi$ est clos à gauche, $\Psi$ clos à droite.
On va voir plus bas \inserted{(voir \unsure{page|pag.} 31)} que cela implique (donc
équivaut à)
\[
\Phi=\Psi^*,\qquad\Psi=\Phi_*.
\tag*{(4.1.3.)}
\]
\scanpage{155}
\mspage{29}
% archive sheet number: 154

Un précouple de Quillen clos est aussi appelé un
couple de Quillen.

\uline{Proposition} 4.2. Soit $(\Phi,\Psi)$ un précouple de Quillen
dans $M$, et soit $j\in\mathrm{Fl}\,M$.$^{*}$ On a alors les conditions
équivalentes sur $j$ :
% slanted
\margin{NB On utilise seulement que $\Phi\longleftrightarrow\Psi$, et que $j$ puisse se mettre sous la forme $pi$, avec $i\in\Phi$, $p\in\Psi$.}

a) $j\in\widetilde\Phi$

b) $j\in\Psi^*$

c) $j$ est facteur direct d'un $i\in\Phi$

c$'$) $j$ est facteur direct d'un $i\in\Phi$ ayant
même source que $j$, plus précisément il
existe un diagramme

\diagram{\begin{tikzcd}[ampersand replacement=\&,column sep=large]
\& A \arrow[r,"{\unsure{\Delta|j}}"] \arrow[d,"\mathrm{id}_A"'] \& B \arrow[d,"q"']\\
\text{(4.2.1.)} \& A \arrow[u,bend right=36,dashed,"\mathrm{id}_A"'] \arrow[r,"i\unsure{{}^{*}}"] \& B' \arrow[u,bend right=30,dashed,"p"']
\end{tikzcd}}
avec $i\in\Phi$, et
\[
qj=i,\quad pi=j,\quad pq=\mathrm{id}_B,
\tag*{(4.2.\unsure{2|1})}
\]
de sorte que $j$ est un facteur direct de $i\in\Phi$,
par $(\mathrm{id}_A,q)$ admettant la rétraction $(\mathrm{id}_A,p)$.

d) Pour tout diagramme commutatif traits pleins
\diagram{\begin{tikzcd}[ampersand replacement=\&,column sep=small]
\& A \arrow[rr,"j"] \arrow[dr,"i"'] \& \& B \arrow[dl,bend left=65,dotted,"q"]\\
\text{(4.2.3.)} \& \& B' \arrow[ur,"p"'] \&
\end{tikzcd}}
i.e. tte factorisation de $j$ en $pi$, si on a
$p\in\Psi$, il existe une section $q$ de $B'$ sur $B$ \inserted{dans $\unsure{{}_{A\backslash}M|A\backslash M}$},
i.e une flèche $q:B\longrightarrow B'$ telle que
\[
pq=\mathrm{id}_B\quad\text{et}\quad qj=i.
\]

Démonstration de la proposition. Elle se fait par la
suite d'implications suivantes, toutes ($\pm$) tautologiques :
\[
c'\Longrightarrow c\Longrightarrow a\Longrightarrow b\Longrightarrow d\Longrightarrow c',
\]
\scanpage{156}
\mspage{30}
% archive sheet number: 155

dont les \unsure{quatre} premières utilisent seulement
l'hypothèse $\Phi\longleftrightarrow\Psi$ (de a) de la définition
4.1.), et la dernière seulement, $d\Rightarrow c'$, fait appel
à l'hypothèse de factorisation de Quillen (b) de
4.1.) ; et encore il suffit de la supposer pour $j$ seulement !.

$c')\Rightarrow c)$ tautologie pure, $c)\Rightarrow a)$ car on sait
que $\widetilde\Phi$ est stable par facteurs directs (et contient $\Phi$),
$a\Rightarrow b$ car $\widetilde\Phi\subset\Psi^*$ ($\widetilde\Phi$ étant par définition la plus
petite partie close à g. de $\mathrm{Fl}\,M$ contenant $\Phi$,
et $\Psi^*$ étant close à gauche), $b\Rightarrow d$ car on écrit
le diagramme donné 4.2.2. (traits pleins)
comme

\diagram{\begin{tikzcd}[ampersand replacement=\&,column sep=large]
\& A \arrow[r,"\unsure{u}"] \arrow[d,"j"',"\in\Psi^*"] \& X \arrow[d,"p\in\Psi"]\\
\text{($\unsure{\alpha|a|\sigma}$)} \& B \arrow[r,"v"'] \arrow[ur,dashed,"q"'] \& Y
\end{tikzcd}}
en y faisant \deleted{\unsure{$B=X$}} $Y=B$, $v=\mathrm{id}_B$ \inserted{et en rebaptisant $X$ \deleted{\unsure{de}} par $B'$, et $u$ par $i$ :} ;
l'hypothèse b) $j\in\Psi^*$ signifie que dans un diagramme \inserted{traits pleins} ($\unsure{\alpha|a}$), il
existe tjrs $q$ telle que $qj=u$, $pq=v$,
et en l'exprimant dans le cas particulier
où $v=\mathrm{id}_B$, on trouve d).

Reste à prouver $d)\Rightarrow c')$, en utilisant
cette fois l'hypothèse de factorisation sur $j$, i.e.
qu'on peut écrire $j$ sous la forme $pi$ :
\diagram{\begin{tikzcd}[ampersand replacement=\&,column sep=small]
 A \arrow[rr,"{\unsure{\Delta|j}}"] \arrow[dr,"i"'] \& \& B \arrow[dl,bend left=65,dashed,"q"]\\
 \& B' \arrow[ur,"p"'] \& {,}
\end{tikzcd}}
avec $i\in\Phi$, $p\in\Psi$. En vertu de
d), \inserted{comme $p\in\Psi$} il existe donc $q:B\longrightarrow B'$ satisfaisant
\scanpage{157}
\mspage{31}
% archive sheet number: 156

\[
pq=\mathrm{id}_B,\quad qj=i,
\]
ce qui prouve c$'$), puisque $i\in\Phi$.

\uline{Corollaire} 1: Soit $(\Phi,\Psi)$ un pré-couple de Quillen.
Alors on a
\[
\begin{aligned}
\widetilde{\Phi}=\Psi^*&=\text{ens. des objets de }\underline{\mathrm{Fl}}\,M\text{ qui sont}\\
&\phantom{=}\text{facteurs directs d'éléments de }\Phi,
\end{aligned}
\tag*{(4.3.)}
\]
et dualement
\[
\begin{aligned}
\overline{\Psi}=\Phi_*&=\text{ens. des objets de }\underline{\mathrm{Fl}}\,M\text{ qui sont}\\
&\phantom{=}\text{facteurs directs d'éléments de }\Psi.
\end{aligned}
\tag*{(4.4)}
\]

\uline{Corollaire 2} : Les conditions suivantes sur le
précouple de Quillen $(\Phi,\Psi)$ sont équivalentes

a) Le précouple est \uline{clos} i.e \deleted{c'est un couple de Quillen.} \inserted{$\Phi=\widetilde\Phi$, $\Psi=\overline\Psi$}

b) $\Phi=\Psi^*$, $\Psi=\Phi_*$

c) $\Phi$ et $\Psi$ sont stables par facteurs directs.

Soit maintenant
\[
W\subset\mathrm{Fl}\,M
\tag*{(4.5.)}
\]
un localiseur, satisfaisant les conditions habituelles
\[
\left\{\begin{aligned}
&\text{a) }W\text{ contient les iso}\\
&\text{b) si }x\overset f{\unsure{-|\longrightarrow}}y\overset g{\unsure{-|\longrightarrow}}z\text{ dans }M,\text{ si }\deleted{\unsure{gf}}\text{ deux parmi }f,g,gf\\
&\qquad\text{sont \unsure{des|dans} }W,\text{ alors \unsure{la|le} troisième aussi}
\end{aligned}\right.
\tag*{(4.6.)}
\]
($W$ modérément saturée) \deleted{Rappelons} Supposons
de plus \unsure{donnés|données} \deleted{deux} \inserted{quatre autres} parties
\[
C,TC,\quad F,TF
\tag*{(4.7)}
\]
de $\mathrm{Fl}\,M$, satisfaisant les conditions suivantes
\[
\left\{\begin{aligned}
&\deleted{TC\subset C,\quad TF\subset F}\\
&\text{a) }(C,TF)\text{ et }(\unsure{TC|T},F)\text{ sont des précouples de Quillen}\\
&\text{b) }TC\subset C\cap W,\quad TF\subset F\cap W\\
&\deleted{\unsure{\text{c) si }p\in F\cap W}\text{ \unsure{et si} }p=\unsure{ji|pi},\text{ avec}}\\
&\qquad\deleted{\unsure{i\in TC,\ p\in TF}}.
\end{aligned}\right.
\tag*{(4.8.)}
\]
\scanpage{158}
\mspage{32}
% archive sheet number: 157

Ces propriétés s'explicitent ainsi :

\uline{M1} Propriétés de relèvement-prolongement
\[
C\unsure{\longleftarrow|\longleftrightarrow}TF\quad\text{et}\quad TC\longleftrightarrow F
\]

\uline{M2} Propriété de factorisation de \deleted{toute} flèche
$f$ de $M$ en $f=pi$, avec, au choix, $i\in C,p\in TF$
ou $i\in TC,\ p\in F$.

\[
M'\qquad TC\subset C\cap W,\quad TF\subset F\cap W
\]
\deleted{\uline{Ex} $W=TF\circ TC$ (\unsure{axiome de saturation})}

C'est donc une partie des axiomes d'une
str. de Quillen faible \deleted{\unsure{aussi}} pour $W$
($\overline{\underline{\mathrm{XIII}}}$ p. 27 déf. 1) -- il manque les axiomes
de stabilité M3, M4, et les axiomes de
``propriété'' (\unsure{propriety}) M$''$, savoir $F\subset\mathrm{Fib}_W$, $C\subset\mathrm{Cof}_W$.

\uline{Proposition 4.3.} Soient $M$ une catégorie, $(C,TF)$ et
$(TC,F)$ deux \inserted{pré-}couples de Quillen dans $M$, satisfaisant
\[
TC\subset C,\quad TF\subset F
\tag*{(4.3.1.)}
\]
et soit $W$ une partie de $\mathrm{Fl}(M)$ dont on ne
suppose pour l'instant que 4.6. b) *, et
% margin note written sideways
\margin{* On utilise seulement que si $x\overset f{\unsure{-|\to}}y\overset g{\unsure{-|\to}}z$, et si $gf\in W$, alors $f\in W\Longleftrightarrow g\in W$. On n'a pas besoin que $W$ contienne les iso, ni qu'il soit stable par composition.}
\[
W\supset TC,TF,\quad\text{donc }\deleted{\unsure{F}}\text{on a}
\tag*{(4.3.2)}
\]
\[
TC\subset C\cap W,\quad TF\subset F\cap W.
\tag*{(4.3.3)}
\]
\deleted{C'est l'axiome de propriété de saturation pour $W$).}
Considérons les conditions de \deleted{factorisation} \inserted{prolongement-relèvement} suivantes
a), a$'$), et les conditions de factorisation b) b$'$) \deleted{\unsure{b}}c), d)
(en se rappelant que $\widetilde{TC}=F^*$, $\overline{TF}=C_*$, en
vertu de prop. 4.2 (cor. 1) appliqué aux \inserted{pré-}couples de
Quillen $(TC,F)$, $(C,TF)$) :

\[
\begin{aligned}
\text{a)}\quad &\deleted{\unsure{F}}\ C\cap W\subset\widetilde{TC}
&&(\text{donc }TC\subset C\cap W\subset\widetilde{TC})\\
\text{a$'$)}\quad &F\cap W\subset\overline{TF}
&&(\text{donc }TF\subset F\cap W\subset\overline{TF})\\
\text{b)}\quad &W\subset\deleted{\overline{\vphantom{TF}}}TF\circ\widetilde{TC}\\
\text{b$'$)}\quad &W\subset\overline{TF}\circ\deleted{\widetilde{\vphantom{TC}}}TC\\
\text{c)}\quad &W\subset\overline{TF}\circ\widetilde{TC}\\
\text{d)}\quad &W\subset TF\circ TC&&(\text{donc }\inserted{\text{par (4.3.2.)}}\ W=TF\circ TC\text{ si }W\text{ est stable}\\
&&&\text{par composition}
\end{aligned}
\]
% slanted
\margin{\uline{NB} \deleted{\unsure{aux}} a) signifie \unsure{juste|aussi} que $(C\cap W,F)$ est encore un précouple de Quillen, a\unsure{$'$)|)} signifie que $(C,F\cap W)$ en est un -- et alors cette paire de précouples de Quillen satisfait cette fois à d) i.e. $W\unsure{\subset|=}TF'\circ TC'$}
\scanpage{159}
\mspage{33}
% archive sheet number: 158

Alors toutes les conditions sauf la dernière sont
équivalentes, (et elles sont \inserted{donc} impliquées par cette dernière).
\inserted{De plus,} Lorsque $TC$ ou $TF$ sont stables par facteurs directs
(p. ex. si $TC$ est clos à gauche, ou $TF$ clos à droite),
toutes les conditions envisagées sont équivalentes
(puisque par cor. 2 à prop. 4.2, on a alors $TC=\widetilde{TC}$
resp. $TF=\overline{TF}$, donc b$'$) resp b) équivaut alors à
d)).

Démonstration. Elle se fait par le diagramme
d'implications
\diagram{\begin{tikzcd}[ampersand replacement=\&,row sep=small,column sep=small]
\& d \arrow[dl,Rightarrow] \arrow[dr,Rightarrow] \&\\
b \arrow[dr,Rightarrow] \& Q \& b' \arrow[dl,Rightarrow]\\
\& c \arrow[dl,Rightarrow] \arrow[dr,Rightarrow] \&\\
a \arrow[uu,Rightarrow] \& \& a' \arrow[uu,Rightarrow]
\end{tikzcd}},
où les \inserted{quatre} implications \deleted{$d\Rightarrow b$ et $d\Rightarrow b'$} \inserted{du carré supérieur $Q$} sont
tautologiques. Il suffit donc d'établir les
\deleted{deux} implications
\[
\deleted{\unsure{b\Rightarrow}} c\Rightarrow a\Rightarrow b,
\]
les implications $\deleted{\unsure{b'\Rightarrow}}c\Rightarrow a'\Rightarrow b'$ en résultant
par dualité. \deleted{\unsure{Mais le}}

\unsure{On|les} \deleted{\unsure{admettons}} \inserted{supposons} donc c) i.e
\[
W\subset\overline{TF}\circ\widetilde{TC},
\]
et prouvons a) i.e que
\[
C\cap W\subset\widetilde{TC}.
\]
Mais si $j\in C\cap W$, pour prouver que $j\in\widetilde{TC}$, écrivons
\deleted{appliquons prop. 4.2. au \inserted{pré-}couple de Quillen
$(TC,F)$, en factorisant $j$ en $pi$, avec
$i\in TC$, $p\in F$. Comme $TC\in W$, on a $i\in W$, donc
\unsure{comme} $pi=j\in W$)}
\[
j=pi\qquad i\in\widetilde{TC},\ p\in\overline{TF},
\]
ce qui est possible par l'hyp. c) et $j\in W$.
\scanpage{160}
\mspage{34}
% archive sheet number: 159

\diagram{\begin{tikzcd}[ampersand replacement=\&,column sep=large]
 A \arrow[rr,"{\unsure{j|\Delta}\in C}"] \arrow[dr,"i\in\widetilde{TC}","\unsure{t|i}"'] \& \& B \arrow[dl,bend left=65,dashed,"q"]\\
 \& B' \arrow[ur,"p\in\overline{TF}"'] \& {.}
\end{tikzcd}}

Utilisons maintenant l'hypothèse $j\in C$ et $C\longleftrightarrow\overline{TF}$,
on trouve donc \inserted{(prop. \unsure{3.2|4.2}, $a\Rightarrow d$)} une flèche $q:B\to B'$ telle
que $pq=\mathrm{id}_B$, $qj=i$, donc $j$ apparaît comme
facteur direct de $i$. Comme $i\in\widetilde{TC}$ et que
$\widetilde{TC}$ est stable par facteurs directs, on en
conclut $j\in\widetilde{TC}$, qed.

$a\Rightarrow b$. On suppose donc
\[
C\cap W\subset\widetilde{TC},
\tag*{($\unsure{\alpha|a}$)}
\]
et on veut prouver b) i.e
\[
W\subset\deleted{\overline{\vphantom{TF}}}TF\circ\widetilde{TC}.
\]
Soit donc $f\in W$, et factorisons $f$ en $pi$, avec
$i\in\deleted{\unsure{T}}C$, $p\in TF$. (OK à cause du précouple de Quillen
$(C,TF)$). Comme $TF\subset W$, on a $p\in W$, et comme
$pi\in W$ on en conclut $i\in W$ (saturation
modérée de $W$ -- on doit quand même
l'utiliser !). On a donc $i\in C\cap W$, donc
$i\in\widetilde{TC}$ par $\alpha$), d'où $f\in TF\circ\widetilde{TC}$, qed.

\uline{Corollaire 1} Supposons $TC$ stable par facteurs directs.
\margin{$\Big|$}Alors \inserted{toutes} les conditions de la prop. 4.3. \inserted{sont équivalentes, et}
équivalent à la relation
\[
\deleted{\unsure{T}}C\cap W=TC,
\tag*{(4.3.4.)}
\]
\margin{$\Big|$}Dualement, si $TF$ stable par facteurs directs,
les conditions de la prop 4.3. \deleted{\unsure{sont}} ttes équivalentes
et équivalent à :
\[
F\cap W=TF.
\tag*{(4.3.4$'$)}
\]
\scanpage{161}
\mspage{35}
% archive sheet number: 160

En effet, par le cor. 2 de prop 4.2, l'hypothèse sur $TC$
implique que $TC=\widetilde{TC}$, donc comme $TC\subset\deleted{\unsure{T}}C\cap W$,
l'hypothèse $C\cap W\subset\widetilde{TC}$ ($=TC$) équivaut à $C\cap W=TC$.

\uline{Cor. 2} : Supposons $C$ et $W$ stables par facteurs
directs. Alors les conditions équivalentes de
la prop 4.3. équivalent à :
\[
C\cap W=\widetilde{TC}.
\tag*{(4.3.5.)}
\]
Dualement, si $\unsure{F|E},W$ sont stables par facteurs directs,
les conditions équiv. de prop. 4.3. équivalent :
\[
F\cap W=\overline{TF}.
\tag*{(4.3.5$'$)}
\]

En effet, si $C$ et $W$ sont stables par fact. directs,
il en est de $=$ de $C\cap W$, et \deleted{comme} dès
lors si $C\cap W\subset\widetilde{TC}$, comme $C\cap W$ contient $TC$,
\deleted{donc} on conclut $\widetilde{C\cap W}=\widetilde{TC}$. Mais
\deleted{comme il a} Comme $C\cap W\unsure{\longleftarrow|\longleftrightarrow} TC$ et $C\cap W$
stable par facteurs directs, $C\cap W$ est clos à g.
par cor. 2 de prop. 4.\unsure{2|1}, \deleted{\unsure{donc}} $C\cap W=\widetilde{C\cap W}$,
d'où $C\cap W=\widetilde{TC}$, qed.

\uline{Cor. 3} Supposons les conditions équivalentes \inserted{a) à c)} de
4.3. satisfaites, et de plus \uline{$W$ stable par
facteurs directs}.* Alors les conditions préliminaires
% slanted
\margin{* En fait, il suffit de supposer que l'on a $W\supset\widetilde{TC},\overline{TF}$, \unsure{et|i.e.} \deleted{\unsure{que}} que tout facteur direct d'un élément de $W$ qui est soit dans $C$, soit dans $F$, est encore dans $W$. Alors $(\widetilde C,\widetilde{TC},\overline F,\overline{TF},W)$ définissent une structure de Qu. close, et on sait que cela implique qu'en fait, $W$ est lui aussi \unsure{alors} stable par facteurs directs}
(4.3.1) et (4.3.2) restent satisfaites quand on
remplace $(C,TF)$ et $(TC,F)$, par leurs
clôtures \inserted{$(C',TF')={}$}$(\unsure{\overline C|\widetilde C},\overline{TF})$ et \inserted{$(TC',F')={}$}$(\widetilde{TC},\overline F)$,
et les conditions équivalentes de la prop. 4.3 sont \inserted{encore} satisfaites
\deleted{\unsure{autre}} pour ces couples, sous les formes fortes
\[
\left\{\begin{aligned}
&C'\cap W=TC'\ (= \widetilde{TC'})\\
&F'\cap W=TF'\ (= \overline{TF'})\\
&W=TF'\circ TC'.
\end{aligned}\right.
\tag*{(4.3.6)}
\]
\scanpage{162}
\mspage{36}
% archive sheet number: 161

Le triple $(W,\unsure{E'|C'},F')$ est un triple de Quillen clos
(correspondant à une ``closed model category'') de
Quillen, \inserted{du moins} si on suppose $M$ stable par $\underrightarrow{\lim}$ finies et
$\underleftarrow{\lim}$ finies, et que $W$ \deleted{satisfasse} est modérément saturé).

% margin note written sideways
\margin{Il en résulte le \fbox{Cor. 4} \deleted{Pour que $(W,C,F)$ soit un \unsure{triple}, de Quillen clos,} \inserted{Soit $W$ un loc. mod. \unsure{saturé}} \deleted{il faut que $C$ clos à g. et $F$ clos à droite, $(C,TF\overset{\mathrm{déf}}=C_*)$} et \unsure{$(TC\overset{\mathrm{déf}}=C\cap W$} et $(C,TF),(TC,F)$ deux \inserted{pré-}couples de Quillen tels que $TC\subset C\cap W$, $TF\subset F\cap W$. Les cond. suivantes sont équiv. a) $TC=C\cap W$ b) $TF=F\cap W$ c) $W=TF\circ TC$ d) $(W,C,\unsure{B|F})$ est un \deleted{\unsure{triple}} de Quillen clos.}

\uline{Démonstration} Les relations (4.3.1) impliquent
\[
\widetilde{TC}\subset\widetilde C\quad\text{et}\quad\overline{TF}\subset\overline F
\]
i.e.
\[
TC'\subset C',\quad T\unsure{P|F}'\subset F',
\]
i.e. (4.3.1) pour les nouveaux couples de
Quillen. \inserted{La relation} 4.3.2. $W\supset TC,TF$, implique
\[
W\supset\widetilde{TC},\overline{TF},\quad\text{i.e.}\quad W\supset TC',TF',
\]
\deleted{\unsure{en}} en vertu du cor. 2 à la prop. 4.2 et du
fait que $W$ est stable par facteurs directs ;
\deleted{comme} donc on a (4.3.2) pour les nouveaux
couples. Comme par hypothèse on a
\[
W=\overline{TF}\circ\widetilde{TC}
\]
et ceci s'écrit
\[
W=TF'\circ TC',
\]
donc les nouveaux couples satisfont aux $\unsure{=\mathrm{s}|\Rightarrow}$ conditions,
qui cette fois, les couples étant clos, prennent
la forme forte (4.3.6).

Prouvons que le triple $(W,C',F')$ satisfait
aux axiomes de Quillen d'une \deleted{structure} \inserted{triple} de
Quillen clos. On a bien, par 4.3.6.
\[
\left\{\begin{aligned}
&C'=\inserted{(TF')^*={}}(F'\cap W)^*,\quad F'=(TC')_{\unsure{*|1}}=(C'\cap W)_*\\
&W=TF'\circ TC'=C'_*\circ {F'}^{*}
\end{aligned}\right.
\]
qui correspond à la définition des ``catégories de
modèles de Quillen \uline{clos}''. Les axiomes
M1, \deleted{\unsure{(pro.}} M2, M3, M4 de Quillen sont \unsure{réalisés|rédigés}
de façon triviale, \deleted{M5 \unsure{par hypoth}} (M3 et M4 puisque
\scanpage{163}
\mspage{37}
% archive sheet number: 162

$C$ et $F$, $TC$ et $TF$ sont clos à g et resp. à dr.). L'axiome
M5 est la saturation \deleted{\unsure{non}} modérée de $W$.

4.\deleted{\unsure{3}}4. Voici l'idée pour construire $\pm$ ``toutes'' les
\deleted{\unsure{couples}} structures de Quillen closes
$(C,TC,F,TF)$ \unsure{associée|associées} à un localiseur donné,
$W\subset M$. On suppose que $M$ satisfait aux
conditions du § précédent, concernant l'existence
de $\underrightarrow{\lim}$ filtrantes, d'une filtration \unsure{cardinale|ordinale},
l'accessibilité de $M$ \deleted{\unsure{(pour)}} \inserted{\unsure{celle|avec}} des éléments de
$\mathrm{Filt}^{\pi_0}$). On part d'une partie petite
\[
C_0\subset\mathrm{Fl}\,M
\]
telle que
\[
C_{0*}\subset W,\tag*{(\unsure{4|5}.4.1)}
\]
et on pose
\[
\left\{\begin{aligned}
&C=\widetilde C_0,\qquad TC=C\cap W.\\
&TF=C_*=C_{0*},\qquad F=(TC)_*.
\end{aligned}\right.\tag*{(4.4.2.)}
\]
\deleted{\unsure{On}}

Il y a \deleted{deux} \inserted{un} problème\unsure{s} seulement, vu
que l'on sait déjà que $(C,TF)$ est un
couple de Quillen par le théorème de factorisation
du § 1 (p. 10) : Il faut la factorisation pour
le couple $(TC,F=(TC)_*)$. À ce moment là,
les conditions préliminaires \inserted{(4.3.1), et (4.3.2) de prop.}
\deleted{\unsure{de}} 4.3. sont
satisfaites
\[
TC\subset C,\quad TF\subset F,\quad W\supset TC,TF\deleted{\unsure{\subset W}},
\]
ainsi que la condition a \quad $C\cap W\subset\widetilde{TC}$, puisque
$C\cap W=TC$. Donc on trouve, par l'équivalence
des conditions a, a', c) et vu que $TC$ et $TF$ sont clos
\[
\unsure{TF|TP}\unsure{\supseteq|=} F\cap W,\qquad W=TF\circ TC
\]
